\قسمت{محدوده‌ها}

مهدی به تازگی با بحث محدوده‌ها\پانویس{scope}
در درس برنامه‌نویسی پیشرفته آشنا شده است.
مهدی معتقد است که بین جاوا و جاوا اسکریپت تفاوت زیادی وجود ندارد زیرا نام‌های آن‌ها بسیار مشابه است.
او با دیدن قطعه کد زیر دچار مشکلات روحی شده است.
به او توضیح دهید که چرا خروجی کد با آنچه در جاوا دیده است متفاوت است
و کد را نیز به گونه‌ای تغییر دهید تا مقصود مهدی را برساند.

\begin{latin}
\begin{minted}[]{js}
function hello() {
  var i = 10
  if (i == 10) {
     var i = 20
  }
  console.log(i)
}
\end{minted}
\end{latin}

\نمره{۳}

% پاسخ پیشنهادی است و پیش از استفاده در تصحیح نیاز به بازبینی دارد.
\begin{پاسخ}

خروجی \متن‌لاتین{20} است.
در جاوا اسکریپت \متن‌لاتین{var} محدوده‌ی تابعی دارد نه بلوکی؛ اعلان درون \متن‌لاتین{if} به بالای تابع منتقل می‌شود\پانویس{Hoisting} و همان متغیر بیرونی را مقداردهی می‌کند.
در جاوا، متغیر داخل بلوک متغیری جداست و با پایان بلوک از بین می‌رود.

برای رسیدن به مقصود مهدی از \متن‌لاتین{let} استفاده می‌کنیم که محدوده‌ی بلوکی دارد:

\begin{latin}
\begin{minted}[]{js}
function hello() {
  let i = 10;
  if (i == 10) {
    let i = 20;
  }
  console.log(i); // 10
}
\end{minted}
\end{latin}

\شروع{فقرات}
\فقره توضیح محدوده‌ی تابعی \متن‌لاتین{var}: ۵۰ درصد نمره
\فقره اصلاح کد با \متن‌لاتین{let}: ۵۰ درصد نمره
\پایان{فقرات}

\end{پاسخ}
